\documentclass[../main.tex]{subfiles}

\begin{document}

\section{Conclusiones y Líneas Futuras}

\subsection{Conclusiones}

El presente Trabajo de Fin de Grado ha permitido diseñar, implementar y validar un prototipo funcional de casco inteligente orientado a la monitorización de impactos craneales en deportes de combate.

El desarrollo del proyecto ha implicado la integración de conocimientos de biomecánica, electrónica, programación de sistemas embebidos, comunicación inalámbrica y diseño mecánico, abordando el problema desde una perspectiva multidisciplinar propia de la Ingeniería de la Salud. La transición desde el marco teórico de la lesión cerebral traumática, donde la aceleración rotacional se identifica como variable biomecánica relevante, hasta una implementación embebida real ha supuesto uno de los principales retos técnicos del trabajo.

Desde el punto de vista técnico, se han alcanzado los objetivos planteados inicialmente:

\begin{itemize}
    \item Selección e integración justificada de sensores inerciales adecuados para la medición de aceleraciones y velocidades angulares.
    \item Desarrollo de firmware capaz de detectar eventos de impacto, estimar aceleración rotacional mediante derivación discreta y clasificar impactos según rangos biomecánicos definidos a partir de la literatura científica.
    \item Implementación de comunicación inalámbrica BLE para transmisión estructurada de eventos y resúmenes de sesión.
    \item Diseño y fabricación de un encapsulado funcional mediante modelado CAD e impresión 3D, resolviendo problemas reales de integración mecánica.
    \item Validación experimental en entorno real de entrenamiento, permitiendo evaluar el comportamiento del sistema bajo condiciones dinámicas no idealizadas.
\end{itemize}

Los resultados obtenidos muestran que, en condiciones de fijación mecánica adecuada, el sistema detecta impactos de forma consistente y clasifica mayoritariamente los eventos en zonas de baja severidad durante entrenamiento habitual. Asimismo, las pruebas realizadas han puesto de manifiesto la importancia crítica del acoplamiento mecánico entre sensor y casco, confirmando que la fiabilidad de la medición depende directamente de la estabilidad estructural del conjunto, en línea con lo descrito en la literatura sobre dispositivos instrumentados \cite{OConnor2017}.

La interpretación de eventos clasificados en zonas de mayor severidad sin correlato clínico ha permitido abordar con criterio crítico la relación entre métricas cinemáticas y lesión, reforzando la evidencia científica de que no existe un umbral determinista único de conmoción cerebral \cite{Guskiewicz2011}.

El proyecto no pretende constituir una herramienta clínica definitiva, sino establecer una base tecnológica coherente para la cuantificación objetiva del movimiento craneal en contextos deportivos.

En definitiva, el trabajo realizado evidencia que es posible trasladar fundamentos biomecánicos complejos a una solución embebida funcional, identificando tanto su potencial como sus limitaciones reales dentro de un marco riguroso de ingeniería aplicada.

\subsection{Líneas Futuras}

El sistema desarrollado abre diversas posibilidades de mejora y evolución tecnológica:

\begin{itemize}
    \item Evaluar plataformas de procesamiento con mayor capacidad de muestreo y memoria, que permitan registrar señales con mayor resolución temporal y ampliar el análisis interno.
    
    \item Explorar configuraciones alternativas con mejor acoplamiento mecánico al cráneo, como dispositivos integrados en protectores bucales instrumentados, descritos en la literatura como soluciones con menor error de desacoplamiento.
    
    \item Incorporar almacenamiento interno que permita conservar la señal completa de cada evento para análisis posterior y estudios longitudinales.
    
    \item Implementar perfiles individuales de deportista y considerar el historial acumulado de impactos en la generación de recomendaciones, permitiendo umbrales adaptativos basados en exposición previa.
    
    \item Desarrollar una aplicación externa específica para visualización avanzada, análisis longitudinal y estudio de patrones de exposición acumulativa.
    
    \item Investigar métricas adicionales derivadas de la señal rotacional como herramienta de estudio biomecánico, ampliando el enfoque actual basado en el valor pico.
    
    \item Optimizar el encapsulado mecánico mediante rediseño estructural que garantice fijación robusta y repetible, minimizando el efecto del desacoplamiento observado durante la validación.
\end{itemize}

Estas líneas futuras permitirían evolucionar el prototipo actual hacia un sistema más robusto, preciso y orientado tanto a la investigación biomecánica como a aplicaciones prácticas en el ámbito deportivo, manteniendo siempre una perspectiva preventiva y basada en evidencia científica.

En última instancia, el presente trabajo demuestra que la ingeniería aplicada puede transformar fundamentos biomecánicos complejos en herramientas tecnológicas reales, aportando objetividad y criterio cuantitativo a la práctica deportiva.


\end{document}